\documentclass[11pt]{article}
\usepackage{amsmath,amssymb,amsthm,hyperref}
\newtheorem{proposition}{Proposition}
\newtheorem{lemma}[proposition]{Lemma}
\newcommand{\C}{\mathbb C}
\newcommand{\Ga}{\mathbb G_a}
\newcommand{\Gm}{\mathbb G_m}
\begin{document}
\title{Cross-programme deductions: filtered markers, Gaussian image dimension, and torus targets}
\author{Internal research audit}
\date{15 November 2026}
\maketitle
These are written internal deductions from the specified repository inputs.
They are not entries in the mathematical claim registry and have not been
independently reviewed or formally verified. Finite algebra checks below
are transcription controls, not proofs of universal or topological statements.
Priority has not been established. The shared MATH\_STATUS.json remains the
mathematical-status authority.

\section{The missing conversion from curve degree to matrix degree}
Give $R=\C[x,y,z,w]/(xw-yz-1)$ its minimum representative total-degree
filtration. For a finite right subgroup $\Gamma\subset SL_2$ give
$A=R^\Gamma$ the induced filtration.
\begin{lemma}
The associated graded rings of $R$ and $A$ are domains. Consequently, for
nonzero $f,g\in A$, $\deg(fg)=\deg f+\deg g$. If $h$ is nonconstant and
$p\in\C[T]$ has degree $d>0$, then $\deg p(h)=d\deg h$.
\end{lemma}
\begin{proof}
The highest part of $(xw-yz-1)b$ is $(xw-yz)b_{\rm top}$. Thus
$\operatorname{gr}R=\C[x,y,z,w]/(xw-yz)$, a domain. The group action preserves
the filtration. Averaging over the finite group makes invariants exact
on each filtration quotient, so
$\operatorname{gr}(R^\Gamma)=(\operatorname{gr}R)^\Gamma$, also a domain.
The highest product symbol is nonzero. In $p(h)$ its top power has
strictly greater degree than every other power and cannot cancel.
\end{proof}

The existing normalization theorem in
\texttt{COMPLEX\_HEIGHT\_IMAGES\_AND\_CAUCHY\_RIGIDITY.md} says that a
nonconstant phase on either full exceptional quotient with one-dimensional
compact image is $P=p(h)$ for a regular real-valued seed $h$. Its finite
curve-marker theorem says that a nonzero pure such phase has an infinite
nonzero marked tail for one fixed $h^j$, $1\le j<d=\deg p$.
The existing Cauchy-transform exclusion says its image $p(h(M))$ must
separate the complex plane.

\begin{proposition}
Let $P\ne0$ be pure-null on $SL_2/2O$ or $SL_2/2I$, and suppose its compact
image has real dimension at most one. Put $D=\deg P$ in matrix entries.
Then $D\ge24$ for $2O$, and $D\ge36$ for $2I$. A fixed bad marker $Q$ exists
with $\deg Q\le D-8$, respectively $\deg Q\le D-12$.
At the first possible degree the normalized polynomial $p$ is cubic and
$\deg h$ is exactly $8$, respectively $12$; a bad marker is $h$ or $h^2$.
\end{proposition}
\begin{proof}
Write $P=p(h)$ and $H=\deg h$. The existing exact character/Molien input
GEXCEPTIONALDEGREES20261115 gives $H\ge8$, respectively $H\ge12$.
A polynomial of degree at most two maps a real interval to a segment or
an injective arc: for a quadratic,
$p(t)-p(s)=(t-s)(a(t+s)+b)$. If $a,b$ are real-linearly independent it is
injective; otherwise its image lies in an affine real line. Such an image
does not separate the plane. The retained Cauchy exclusion therefore
forces $d\ge3$. The lemma gives $D=dH\ge3H$. For the existing fixed bad
marker $h^j$, $\deg h^j=jH\le(d-1)H=D-H$. Equality of $D$ with $24$ or $36$
forces $d=3$ and the minimum $H$.
\end{proof}
These bounds concern thin images of the PHASE, not of its auxiliary Hopf
height, and do not exclude full-image counterexamples in smaller degrees.
The thresholds are necessary, not existence statements.

\section{A real-curve Gaussian phase has a small bad marker}
Use the ordinary positive real Gaussian functional on $\C[X_1,\ldots,X_n]$.
\begin{proposition}
If $P\ne0$ has all positive pure Gaussian moments zero and its image on
$\mathbb R^n$ has real dimension at most one, it is not eventually safe.
There is a real polynomial $h$ and a polynomial $p$ of degree $d\ge2$ with
$P=p(h)$, and one fixed marker $h^j$, $1\le j<d$, has infinitely many nonzero
marked moments. Its total degree is at most $\deg P-1$.
\end{proposition}
\begin{proof}
Conjugation fixes real coordinates. Rank at most one gives algebraic
dependence of $P$ and $\overline P$ by the characteristic-zero differential
criterion and real Zariski density. Normalize $\C[P,\overline P]$ in its
fraction field. It embeds in the normal polynomial ring; its curve is
unirational, and its units are constant. Its completion is $\mathbb P^1$
and exactly one point is deleted, so the normalization is $\C[h_0]$.
The conjugation-stable affine coordinate can be changed to a real polynomial
$h$, by the same affine involution calculation as the retained compact
normalization proof. This algebraic part uses no compactness.

The Gaussian pairing on $\C[h]$ is nondegenerate, because
$\mathbb E[q(h)\overline q(h)]>0$ for every nonzero polynomial $q$.
All these polynomial integrals are finite. This ring is free of rank $d$
over $\C[p(h)]$, with basis $1,h,\ldots,h^{d-1}$. The retained finite-module
marker lemma GFINITEMODULEMARKER20261115 would force the nonzero $P$ to
be nilpotent if all basis tails eventually vanished. Purity removes the
marker $1$, giving the specified $j$. Ordinary polynomial degrees multiply
under composition: $\deg P=d\deg h$. Thus $j\deg h\le\deg P-\deg h\le\deg P-1$.
\end{proof}
In particular every nonzero pure ternary Gaussian cubic has image with
interior, by G3CUBICALL20261115. The same conclusion holds in the admitted
Gaussian shear orbits of such cubics, despite their unbounded degree.
This is a geometric exclusion and marker bound, not a construction of a
new pure phase or a solution of general GMC.

\section{All torus-quotient-valued maps factor through the old signal}
Put $U=\operatorname{Spec}\C[S,Y,(1-SY)^{-1}]$ and $r=1-SY$.
Here a torus quotient means $Z=T/F$ with $T$ a complex algebraic torus and
$F$ a finite group acting freely by algebraic variety automorphisms,
including translations, exactly as in the torus receiver programme.
\begin{proposition}
For every $a\ge0$ every regular map $f:U\times\mathbb A^a\longrightarrow Z$
factors uniquely through $r:U\times\mathbb A^a\longrightarrow\Gm$.
In particular its image has dimension at most one, so it cannot dominate
such a target of dimension at least two. No etaleness, finiteness, or
cofinite-image hypothesis on $f$ is required.
\end{proposition}
\begin{proof}
The existing UAPRODUCTTOPO1 computation gives $\pi_1(U)=\mathbb Z$ with
$r_*$ an isomorphism onto $\pi_1(\Gm)$; affine factors do not change it.
Every connected finite etale degree-$m$ cover is consequently
\[
 U_m\times\mathbb A^a,\qquad
 U_m=\operatorname{Spec}\C[S,Y,v,v^{-1}]/(SY-1+v^m),\quad r=v^m.
\]
These explicit covers are connected, smooth and finite etale, and realize
all finite-index subgroups $m\mathbb Z$. The usual full faithfulness from
finite etale complex covers to topological covers identifies them
algebraically: the graph of a cover isomorphism is a connected component
of the finite etale fibre product, hence an algebraic open-and-closed
component.

Their units are precisely $c v^b$. Indeed localization at $S$ is the
Laurent ring $\C[S^{\pm1},v^{\pm1}]$, whose units are $cS^a v^b$.
Each prime $(S,v-\zeta)$, $\zeta^m=1$, has $S$ of positive order and $v$ a
unit. An element and its inverse can both be regular there only if $a=0$.
Adding polynomial affine variables gives no new units.

Pull back $T\to Z$ along $f$ and take a connected component. It is one of
these covers. The lifted map to $T$ is specified by invertible coordinate
functions, all of the form $cv^b$, so factors through $v:U_m\to\Gm$.
Its composite to $Z$ is invariant under the deck transformations
$v\mapsto\zeta v$, with $S,Y$ and the affine variables fixed. Since $v$ is
surjective, the resulting map $\Gm\to Z$ is $\mu_m$-invariant and factors
through $v\mapsto v^m$. Descent gives $f=\psi\circ r$. Surjectivity of $r$
proves uniqueness.
\end{proof}
This yields a whole-map obstruction even for unit-free torus quotients.
For the proper native intermediate surface this overlaps with the existing
negative logarithmic-dimension obstruction; it is not an additional proof
of JC(2) or canonical atomicity. It does not classify arbitrary surfaces
with virtually abelian fundamental group.

\section{A stronger ring statement at the octahedral lowest type}
The previously identified radial transfer is now registered as
GEXCEPTIONALRADIAL20261115/GEXCEPTIONALUNIFORM20261115 in
\texttt{GAUSSIAN\_RADIAL\_TRANSFER\_AND\_EXCEPTIONAL\_UNIFORMITY.md}.
It is not counted as a fresh discovery in this audit.

\begin{proposition}
On the full octahedral degree-at-most-eleven phase space, every all-order
pure-null phase over an arbitrary complex coefficient algebra is eventually
safe, and one common integer linear cutoff works for all such algebras.
The full ordinary mixed-cutoff ideals eventually equal the full ordinary
pure ideal, rather than only having the same radical.
\end{proposition}
\begin{proof}
GOCTAHEDRALPREFIX20261115 identifies the space as
$\C\oplus V_8$, with every field-pure point safe and common field slope
$1/2$. GOCTAHEDRALIDEAL20261115 gives the full ordinary pure ideal with
least prefix ten. Transport this parameter family through the registered
radial identity into four Gaussian variables. Prescribed cutoff ideals,
marker degrees and the full ordinary pure ideal agree there.
Apply the positive-family theorem MVPOSITIVEIDEAL20261115 to that
Gaussian coefficient family: on the pure quotient the radical safe ideal
is the nilradical, with a fixed finite exponent, so the complete ordinary
cutoff ideals eventually equal the pure ideal. Transfer back.
\end{proof}
This does not retain the sharp field slope $1/2$ over nonreduced algebras;
no new numerical all-ring cutoff or nilpotence exponent is computed.
It strengthens an already positive degree range, without advancing the
unknown field case at degree twelve or settling the full exceptional model.

\section{Audit boundary and suggested use}
The phase-degree $24/36$ exclusion and marker-degree conversion are new
cross-programme deductions to submit to the usual source/registry review.
The Gaussian thin-image statement adds a bounded marker criterion and
excludes that geometry inside all already positive classes. The torus
factorization is a general whole-source lemma with an explicitly identified
native overlap. The octahedral ring consequence is a synthesis of registered
positive results. No entry in MATH\_STATUS.json or active agent proof is
modified by this audit. The numerical checks in the companion Python file
do not formalize these arguments or establish novelty.
\end{document}
